% =====================================================================
%  INTRODUCCIÓN — capítulo sin número
% =====================================================================
\chapter*{Introducción}
\addcontentsline{toc}{chapter}{Introducción}  % Lo agrega al índice
\markboth{Introducción}{Introducción}
% Sin \label: un capítulo sin número no se puede citar con \cref
% (tomaría el número del capítulo anterior). Escriba "la introducción".

% Contenido sugerido: planteamiento del problema, justificación, pregunta
% de investigación, objetivos, hipótesis y organización del documento.
% En un capítulo sin número use \section* (secciones sin número).

\section*{Planteamiento del problema}

\lipsum[1]

\section*{Objetivos}

\textbf{Objetivo general.} \lipsum[2][1-3]

\textbf{Objetivos específicos:}
\begin{enumerate}
  \item \lipsum[3][1]
  \item \lipsum[3][2]
  \item \lipsum[3][3]
\end{enumerate}

\section*{Hipótesis}

\lipsum[4][1-4]

\section*{Estructura de la tesis}

% Con \cref los números se actualizan solos si cambia el orden.
La tesis se organiza en cuatro capítulos. El \cref{cap:capitulo1} \ldots; el
\cref{cap:capitulo2} \ldots; el \cref{cap:capitulo3} \ldots y el
\cref{cap:capitulo4} \ldots. Finalmente se presentan las conclusiones y
recomendaciones.
